% ============================================================
% Script de soutenance — Myriam (diapos 1 à 4)
% VERSION ADAPTÉE DYSLEXIE :
%   - police OpenDyslexic (fichiers dans fonts/), corps 12 pt
%   - interligne 1,6 et paragraphes espacés
%   - lignes courtes (marges larges), texte non justifié
%   - pas d'italique ; phrases courtes ; listes plutôt que pavés
% Compilation : xelatex script_myriam_dyslexie.tex
% (xelatex obligatoire : la police est chargée via fontspec)
% ============================================================
\documentclass[12pt,a4paper]{article}

\usepackage[french]{babel}
\usepackage[margin=3.2cm]{geometry}
\usepackage{amsmath,amssymb}
\usepackage{booktabs}
\usepackage{parskip}
\usepackage{setspace}
\usepackage{fontspec}
\setmainfont{OpenDyslexic}[
  Path           = fonts/,
  UprightFont    = OpenDyslexic-Regular.otf,
  BoldFont       = OpenDyslexic-Bold.otf,
  ItalicFont     = OpenDyslexic-Italic.otf,
  BoldItalicFont = OpenDyslexic-BoldItalic.otf,
]
\usepackage{enumitem}

\setstretch{1.6}
\raggedright
\setlength{\parskip}{0.9em}
\setlist{itemsep=0.5em, topsep=0.5em}

\newcommand{\diapo}[2]{\section*{Diapo #1 --- #2}}
\newenvironment{comprendre}%
  {\par\medskip\noindent\hrulefill\par\nopagebreak\textbf{Pour bien comprendre :}\par}%
  {\par\nopagebreak\noindent\hrulefill\par\medskip}

\title{\textbf{Script de soutenance --- Myriam}\\
\large Effet tunnel et temps de traversée d'une barrière\\
\normalsize Diapos 1 à 4 --- environ 5 minutes}
\author{}
\date{Soutenance juin 2026 --- oraux du 22 au 26 juin}

\begin{document}
\maketitle

\subsection*{Qui dit quoi}

\begin{itemize}
\item \textbf{Myriam : diapos 1 à 4. Le problème, le paquet d'ondes, la barrière. Environ 5 min.}
\item Sheryne : diapos 5 à 8. La méthode numérique et les premiers résultats. Environ 4,5 min.
\item Mattéo : diapos 9 à 12. L'effet Hartman et la conclusion. Environ 5,5 min.
\end{itemize}

Total : 15 minutes, puis 10 minutes de questions.

Les phrases \textbf{en gras} sont les messages essentiels.
Si le temps presse, on peut raccourcir le reste, mais jamais le gras.

% ------------------------------------------------------------
\clearpage
\diapo{1}{Page de titre (30 secondes)}

\textbf{Texte à dire :}

Bonjour à toutes et à tous.

Nous allons vous présenter notre projet numérique de physique moderne.
Il porte sur \textbf{l'effet tunnel et le temps de traversée
d'une barrière de potentiel} par un paquet d'ondes gaussien,
à une dimension.

Je suis Myriam Bensaid.
Je travaille avec Sheryne Ouarghi et Mattéo Knopes.

Le plan :

\begin{itemize}
\item je pose d'abord le problème et les outils théoriques ;
\item Sheryne présente ensuite la méthode numérique et les premiers résultats ;
\item Mattéo termine par l'étude du temps de traversée et la conclusion.
\end{itemize}

% ------------------------------------------------------------
\clearpage
\diapo{2}{Problématique (1 minute)}

\textbf{Texte à dire :}

En mécanique classique, c'est simple.
Une particule d'énergie $E$ arrive sur une barrière de hauteur $V_0$.
Si $V_0 > E$, elle rebondit. Toujours.

En mécanique quantique, c'est différent.
La particule a une probabilité non nulle de passer de l'autre côté.
C'est l'effet tunnel.

Notre question va un cran plus loin :
\textbf{combien de temps met une particule quantique pour traverser
une barrière qu'elle n'a classiquement pas le droit de franchir ?}
Cette question paraît simple.
Elle est encore débattue aujourd'hui.
Et elle réserve des surprises.

Notre démarche, en cinq étapes :

\begin{enumerate}
\item construire un paquet d'ondes gaussien ;
\item déterminer les états stationnaires de la barrière ;
\item résoudre numériquement l'équation de Schrödinger dépendante du temps ;
\item mesurer le temps de traversée et le comparer à la théorie ;
\item étudier l'influence des paramètres de la barrière.
\end{enumerate}

Tout le projet est en unités adimensionnées : $\hbar = m = 1$.
Cela allège les équations sans changer la physique.

\begin{comprendre}
Pourquoi la question du temps est-elle difficile ?
Sous la barrière, l'onde n'oscille pas : elle s'amortit.
On dit qu'elle est « évanescente ».
Il n'y a donc pas de vitesse évidente à l'intérieur.

Il faut donc choisir une définition du temps de traversée.
La nôtre : le temps entre le passage du pic du paquet
à l'entrée de la barrière, et son passage à la sortie.
\end{comprendre}

% ------------------------------------------------------------
\clearpage
\diapo{3}{Le paquet d'ondes gaussien (1,5 minute)}

\textbf{Texte à dire :}

Première étape : représenter une particule.

Une onde plane unique ne convient pas.
Elle est présente partout dans l'espace.
Pour localiser la particule, on superpose une infinité d'ondes planes.
C'est l'intégrale à l'écran.
Chaque onde de vecteur d'onde $k$ a un poids gaussien $g(k)$,
centré sur une valeur moyenne $k_0$.

\textbf{Ce paquet a quatre propriétés clés :}

\begin{itemize}
\item \textbf{La vitesse de groupe} : $v_g = k_0$ dans nos unités.
      C'est la vitesse du centre du paquet.
      Elle joue le rôle de la vitesse de la particule.
\item \textbf{La vitesse de phase} : les oscillations internes
      avancent à la moitié de $v_g$.
\item \textbf{La largeur spectrale} : $\Delta k = 1/(2\sqrt{a_{\rm wp}})$.
      Le paramètre $a_{\rm wp}$ contrôle la largeur spatiale.
      \textbf{Plus le paquet est étroit en position,
      plus il est large en vecteur d'onde.}
      C'est le principe d'indétermination de Heisenberg.
\item \textbf{L'étalement} : le paquet s'élargit au cours du temps.
      La relation $\omega = k^2/2$ n'est pas linéaire,
      donc les composantes rapides prennent de l'avance.
\end{itemize}

Sur la figure de droite : les parties réelle et imaginaire du paquet
à l'instant initial.
Des oscillations sous une enveloppe gaussienne.

\begin{comprendre}
Le point à retenir : $\Delta k$.
Toute la fin de l'exposé repose dessus.

Le paquet n'a pas une seule énergie.
Il a une plage d'énergies, de largeur liée à $\Delta k$.
Or la barrière ne traite pas toutes les composantes pareil.
Ce « filtrage » expliquera les écarts entre nos mesures
et la théorie onde plane.
\end{comprendre}

% ------------------------------------------------------------
\clearpage
\diapo{4}{États stationnaires de la barrière (2 minutes --- diapo indispensable)}

\textbf{Texte à dire :}

Deuxième outil, et le jury l'attend : les états stationnaires.
On cherche les solutions d'énergie $E$ fixée, avec $E < V_0$.
La barrière est rectangulaire : hauteur $V_0$, largeur $a$.

L'espace se découpe en trois zones :

\begin{itemize}
\item avant la barrière : une onde incidente, plus une onde réfléchie
      d'amplitude $R$ ;
\item après la barrière : une onde transmise d'amplitude $T$ ;
\item \textbf{dans la barrière : le vecteur d'onde devient imaginaire.}
      La solution n'oscille plus.
      Elle est exponentielle, en $e^{\pm\kappa x}$,
      avec $\kappa = \sqrt{2m(V_0 - E)}/\hbar$.
      C'est l'onde évanescente : elle ne se propage pas, elle s'amortit.
\end{itemize}

Pour relier les trois zones, on impose la continuité de $\psi$
et de sa dérivée, en $x = 0$ et en $x = a$.
Quatre équations, quatre inconnues.
On en tire l'amplitude de transmission $T(k)$, dans l'encadré.

Deux conséquences :

\begin{itemize}
\item \textbf{$|T|^2$ n'est jamais nul : la transmission est toujours
      possible.} C'est la signature mathématique de l'effet tunnel.
\item En régime opaque ($\kappa a \gg 1$), la transmission décroît
      comme $e^{-2\kappa a}$.
      \textbf{Une sensibilité exponentielle à la largeur et à la hauteur.}
      C'est ce qui rend l'effet tunnel si sélectif.
      Et si utile : par exemple le microscope à effet tunnel.
\end{itemize}

Dernier point : $T$ est un nombre complexe.
Il a un module, mais aussi une phase $\varphi_T$.
\textbf{Cette phase contient l'information sur le temps de traversée.}
Mattéo s'en servira pour la prédiction théorique.

Je laisse maintenant la parole à Sheryne pour la résolution numérique.

\begin{comprendre}
Le raccordement est une question classique de jury.
Il faut savoir le refaire au tableau.

La logique :
dans chaque zone, l'équation de Schrödinger est une équation
différentielle linéaire d'ordre 2, à coefficients constants.
Donc la solution est une combinaison de deux exponentielles.
Les conditions de continuité aux deux interfaces fixent les amplitudes.

Pourquoi $\psi'$ est-elle continue ?
Parce que le potentiel reste fini.
La dérivée ne saute que pour un potentiel en pic de Dirac.
\end{comprendre}

% ------------------------------------------------------------
\clearpage
\section*{Vue d'ensemble (pour répondre sur tout l'exposé)}

Le fil rouge, en quatre points :

\begin{enumerate}
\item But : mesurer le temps qu'un paquet d'ondes met à traverser
      une barrière par effet tunnel.
\item Ma partie donne les deux outils.
      Le paquet gaussien : vitesse $v_g = k_0$, largeur spectrale $\Delta k$.
      Et la transmission $T(k)$ : décroissance exponentielle, phase $\varphi_T$.
\item Sheryne : le solveur Crank--Nicolson est fiable
      (norme conservée, $\tau_0$ retrouvé à 4\,\% près).
      Le paquet se sépare en réfléchi (gros) et transmis (petit).
      La transmission mesurée dépasse la théorie onde plane :
      c'est le filtrage spectral.
\item Mattéo : le résultat principal.
      \textbf{Le paquet tunnel sort plus tôt que le paquet libre :
      $\tau_t < \tau_0$. C'est l'effet Hartman.}
      Puis son comportement face à la barrière : saturation avec $a$,
      diminution avec $V_0$. En perspective : l'influence du paquet
      lui-même ($k_0$, $a_{\rm wp}$), étudiée mais pas montrée aujourd'hui.
\end{enumerate}

\textbf{Les nombres à connaître} (cas de référence : $k_0 = 5$, $E = 12{,}5$,
$V_0 = 15$, $a = 1$, $a_{\rm wp} = 1$) :

\begin{center}
\begin{tabular}{lcc}
\toprule
Grandeur & Numérique & Théorique \\
\midrule
Temps libre $\tau_0$ & 0{,}192 & 0{,}200 \\
Temps tunnel $\tau_t$ & 0{,}112 & 0{,}168 (Hartman) \\
Transmission & $8{,}1\times10^{-2}$ & $2{,}5\times10^{-2}$ \\
Saturation de $\tau_t^{\rm th}$ & --- & $\approx 0{,}179$ \\
\bottomrule
\end{tabular}
\end{center}

% ------------------------------------------------------------
\clearpage
\section*{Questions probables du jury --- réponses courtes}

\textbf{Refaire le raccordement des états stationnaires au tableau.}
C'est MA question (diapo 4).
Trois zones : ondes oscillantes dehors, onde évanescente dedans.
Continuité de $\psi$ et de $\psi'$ en $x=0$ et $x=a$.
Quatre équations, quatre inconnues : on en tire $T(k)$.

\textbf{Pourquoi Crank--Nicolson plutôt qu'Euler explicite ?}
Crank--Nicolson est unitaire : la norme est conservée par construction.
Il est stable quel que soit le pas de temps.
Euler explicite amplifie la norme et diverge.
Coût faible : système tridiagonal, méthode de Thomas.
C'est la partie de Sheryne, mais il faut savoir le dire.

\textbf{Avez-vous étudié $k_0$ et $a_{\rm wp}$ ? Les temps négatifs
violent-ils la causalité ?}
Oui, nous l'avons fait, mais sans le garder dans le diaporama,
par souci de temps (perspective de la conclusion).
Non, pas de violation : le pic transmis n'est pas un objet qui se déplace.
Il est reconstruit à partir des composantes rapides du spectre.
Aucune information ne va trop vite.
C'est un effet de filtrage spectral.

\textbf{Pourquoi la transmission mesurée dépasse-t-elle $|T(k_0)|^2$ ?}
Le paquet n'est pas monochromatique.
La barrière transmet beaucoup mieux les composantes rapides.
Quand le paquet devient quasi monochromatique ($a_{\rm wp}$ grand),
la mesure devrait retomber sur $|T(k_0)|^2$ : une piste creusée
mais pas gardée dans le diaporama.

\textbf{Effet Hartman : que se passe-t-il pour $a \to \infty$ ?}
Le temps théorique sature alors que la distance augmente.
La vitesse apparente semblerait énorme.
Mais la transmission s'effondre exponentiellement,
et la mesure du pic perd son sens bien avant.
Pas de transport d'information plus vite que la lumière.

\end{document}
